\section{Implementasi AES dalam C}

\subsection{Struktur AES dan Komponen Utama}

AES (Advanced Encryption Standard) menggunakan struktur Substitution-Permutation Network (SPN) dengan blok 128 bit dan kunci 128/192/256 bit \cite{fips197}. Setiap round terdiri dari empat operasi:

\begin{enumerate}
  \item \textbf{SubBytes}: Substitusi non-linear menggunakan S-box AES
  \item \textbf{ShiftRows}: Permutasi baris untuk diffusion
  \item \textbf{MixColumns}: Operasi linear dalam $\mathrm{GF}(2^8)$
  \item \textbf{AddRoundKey}: XOR dengan round key
\end{enumerate}

\subsection{Representasi Data AES}

Dalam implementasi C, state AES direpresentasikan sebagai matriks 4x4 bytes:
\begin{itemize}
  \item Input: 16 bytes plaintext
  \item State: Matriks 4x4 bytes $s_{r,c}$ (baris $r$, kolom $c$)
  \item Output: 16 bytes ciphertext
\end{itemize}

\begin{table}[htbp]
  \centering
  \small
  \begin{tabular}{|c|c|c|c|}
    \hline
    $s_{0,0}$ & $s_{0,1}$ & $s_{0,2}$ & $s_{0,3}$ \\
    \hline
    $s_{1,0}$ & $s_{1,1}$ & $s_{1,2}$ & $s_{1,3}$ \\
    \hline
    $s_{2,0}$ & $s_{2,1}$ & $s_{2,2}$ & $s_{2,3}$ \\
    \hline
    $s_{3,0}$ & $s_{3,1}$ & $s_{3,2}$ & $s_{3,3}$ \\
    \hline
  \end{tabular}
  \caption{Representasi state AES 4x4.}
\end{table}

\subsection{S-Box AES}

S-Box AES adalah substitusi non-linear 8-bit yang dirancang untuk resistensi terhadap analisis linear dan differential \cite{daemen2001}. S-Box dihasilkan dari inversi multiplicative dalam $\mathrm{GF}(2^8)$ diikuti dengan transformasi affine.

\begin{lstlisting}[style=cstyle, caption={S-Box AES dalam C}]
static const uint8_t sbox[256] = {
    0x63, 0x7c, 0x77, 0x7b, 0xf2, 0x6b, 0x6f, 0xc5,
    0x30, 0x01, 0x67, 0x2b, 0xfe, 0xd7, 0xab, 0x76,
    0xca, 0x82, 0xc9, 0x7d, 0xfa, 0x59, 0x47, 0xf0,
    0xad, 0xd4, 0xa2, 0xaf, 0x9c, 0xa4, 0x72, 0xc0,
    // ... (256 values total)
    0x70, 0x3f, 0xe6, 0xb1, 0x98, 0x11, 0x69, 0xd9,
    0x8e, 0x94, 0x9b, 0x1e, 0x87, 0xe9, 0xce, 0x55,
    0x28, 0xdf, 0x8c, 0xa1, 0x89, 0x0d, 0xbf, 0xe6
};

void sub_bytes(uint8_t state[4][4]) {
    for (int r = 0; r < 4; r++) {
        for (int c = 0; c < 4; c++) {
            state[r][c] = sbox[state[r][c]];
        }
    }
}
\end{lstlisting}

\subsection{Operasi AES dalam C}

\subsubsection{ShiftRows}

ShiftRows menggeser setiap baris ke kiri dengan jumlah yang berbeda:
\begin{itemize}
  \item Baris 0: tidak digeser
  \item Baris 1: digeser 1 posisi
  \item Baris 2: digeser 2 posisi  
  \item Baris 3: digeser 3 posisi
\end{itemize}

\begin{lstlisting}[style=cstyle, caption={ShiftRows dalam C}]
void shift_rows(uint8_t state[4][4]) {
    uint8_t temp;
    
    // Row 1: shift 1
    temp = state[1][0];
    state[1][0] = state[1][1];
    state[1][1] = state[1][2];
    state[1][2] = state[1][3];
    state[1][3] = temp;
    
    // Row 2: shift 2
    temp = state[2][0];
    state[2][0] = state[2][2];
    state[2][2] = temp;
    temp = state[2][1];
    state[2][1] = state[2][3];
    state[2][3] = temp;
    
    // Row 3: shift 3
    temp = state[3][0];
    state[3][0] = state[3][3];
    state[3][3] = state[3][2];
    state[3][2] = state[3][1];
    state[3][1] = temp;
}
\end{lstlisting}

\subsubsection{MixColumns}

MixColumns melakukan operasi linear dalam $\mathrm{GF}(2^8)$ menggunakan matriks tetap:
\[
\begin{bmatrix}
s'_{0,c} \\ s'_{1,c} \\ s'_{2,c} \\ s'_{3,c}
\end{bmatrix}
=
\begin{bmatrix}
02 & 03 & 01 & 01 \\
01 & 02 & 03 & 01 \\
01 & 01 & 02 & 03 \\
03 & 01 & 01 & 02
\end{bmatrix}
\cdot
\begin{bmatrix}
s_{0,c} \\ s_{1,c} \\ s_{2,c} \\ s_{3,c}
\end{bmatrix}
\]

\begin{lstlisting}[style=cstyle, caption={MixColumns dalam C}]
// Galois Field multiplication
uint8_t gmul(uint8_t a, uint8_t b) {
    uint8_t p = 0;
    for (int counter = 0; counter < 8; counter++) {
        if ((b & 1) != 0) {
            p ^= a;
        }
        uint8_t hi_bit_set = (a & 0x80);
        a <<= 1;
        if (hi_bit_set) {
            a ^= 0x1b; // x^8 + x^4 + x^3 + x + 1
        }
        b >>= 1;
    }
    return p;
}

void mix_columns(uint8_t state[4][4]) {
    uint8_t temp[4];
    
    for (int c = 0; c < 4; c++) {
        temp[0] = gmul(0x02, state[0][c]) ^ gmul(0x03, state[1][c]) ^ 
                 gmul(0x01, state[2][c]) ^ gmul(0x01, state[3][c]);
        temp[1] = gmul(0x01, state[0][c]) ^ gmul(0x02, state[1][c]) ^ 
                 gmul(0x03, state[2][c]) ^ gmul(0x01, state[3][c]);
        temp[2] = gmul(0x01, state[0][c]) ^ gmul(0x01, state[1][c]) ^ 
                 gmul(0x02, state[2][c]) ^ gmul(0x03, state[3][c]);
        temp[3] = gmul(0x03, state[0][c]) ^ gmul(0x01, state[1][c]) ^ 
                 gmul(0x01, state[2][c]) ^ gmul(0x02, state[3][c]);
        
        for (int r = 0; r < 4; r++) {
            state[r][c] = temp[r];
        }
    }
}
\end{lstlisting}

\subsubsection{AddRoundKey}

AddRoundKey melakukan XOR antara state dan round key:
\begin{lstlisting}[style=cstyle, caption={AddRoundKey dalam C}]
void add_round_key(uint8_t state[4][4], uint8_t round_key[4][4]) {
    for (int r = 0; r < 4; r++) {
        for (int c = 0; c < 4; c++) {
            state[r][c] ^= round_key[r][c];
        }
    }
}
\end{lstlisting}

\subsection{Key Schedule AES}

Key schedule menghasilkan 11 round keys dari 128-bit kunci untuk AES-128:
\begin{lstlisting}[style=cstyle, caption={Key Schedule AES-128}]
// Rcon untuk AES-128
static const uint8_t rcon[10] = {
    0x01, 0x02, 0x04, 0x08, 0x10, 0x20, 0x40, 0x80, 0x1b, 0x36
};

// RotWord: rotate 32-bit word left by 8 bits
uint32_t rot_word(uint32_t word) {
    return (word << 8) | (word >> 24);
}

// SubWord: apply S-box to each byte
uint32_t sub_word(uint32_t word) {
    return (sbox[(word >> 24) & 0xff] << 24) |
           (sbox[(word >> 16) & 0xff] << 16) |
           (sbox[(word >> 8) & 0xff] << 8) |
           (sbox[word & 0xff]);
}

void key_expansion(uint8_t key[16], uint8_t round_keys[11][16]) {
    uint32_t *words = (uint32_t*)key;
    uint32_t temp;
    
    // Copy original key as first round key
    memcpy(round_keys[0], key, 16);
    
    // Generate remaining round keys
    for (int i = 4; i < 44; i++) {
        temp = words[i - 1];
        
        if (i % 4 == 0) {
            temp = sub_word(rot_word(temp)) ^ rcon[i/4 - 1];
        }
        
        words[i] = words[i - 4] ^ temp;
        
        // Copy to round keys (every 4 words = 16 bytes)
        if ((i + 1) % 4 == 0) {
            memcpy(round_keys[(i + 1)/4], &words[i - 3], 16);
        }
    }
}
\end{lstlisting}

\subsection{Implementasi AES-128 Lengkap}

\begin{lstlisting}[style=cstyle, caption={AES-128 Enkripsi Lengkap}]
#include <stdint.h>
#include <string.h>

// S-Box, Rcon, gmul, sub_bytes, shift_rows, mix_columns, add_round_key, key_expansion
// ... (functions defined above)

void aes_encrypt(uint8_t plaintext[16], uint8_t ciphertext[16], uint8_t key[16]) {
    uint8_t state[4][4];
    uint8_t round_keys[11][16];
    
    // Key expansion
    key_expansion(key, round_keys);
    
    // Copy plaintext to state
    for (int r = 0; r < 4; r++) {
        for (int c = 0; c < 4; c++) {
            state[r][c] = plaintext[r * 4 + c];
        }
    }
    
    // Initial round
    add_round_key(state, (uint8_t(*)[4])round_keys[0]);
    
    // Main rounds (1-9)
    for (int round = 1; round <= 9; round++) {
        sub_bytes(state);
        shift_rows(state);
        mix_columns(state);
        add_round_key(state, (uint8_t(*)[4])round_keys[round]);
    }
    
    // Final round (no MixColumns)
    sub_bytes(state);
    shift_rows(state);
    add_round_key(state, (uint8_t(*)[4])round_keys[10]);
    
    // Copy state to ciphertext
    for (int r = 0; r < 4; r++) {
        for (int c = 0; c < 4; c++) {
            ciphertext[r * 4 + c] = state[r][c];
        }
    }
}

int main() {
    uint8_t plaintext[16] = {0x00, 0x11, 0x22, 0x33, 0x44, 0x55, 0x66, 0x77,
                            0x88, 0x99, 0xaa, 0xbb, 0xcc, 0xdd, 0xee, 0xff};
    uint8_t key[16] = {0x00, 0x01, 0x02, 0x03, 0x04, 0x05, 0x06, 0x07,
                      0x08, 0x09, 0x0a, 0x0b, 0x0c, 0x0d, 0x0e, 0x0f};
    uint8_t ciphertext[16];
    
    aes_encrypt(plaintext, ciphertext, key);
    
    printf("Plaintext:  ");
    for (int i = 0; i < 16; i++) printf("%02x ", plaintext[i]);
    printf("\nCiphertext: ");
    for (int i = 0; i < 16; i++) printf("%02x ", ciphertext[i]);
    printf("\n");
    
    return 0;
}
\end{lstlisting}

\subsection{Optimasi Implementasi}

Implementasi AES dapat dioptimasi dengan beberapa teknik:

\begin{itemize}
  \item \textbf{Tabel Lookup}: Pre-compute MixColumns dan SubBytes
  \item \textbf{Bit-slicing}: Proses multiple bytes secara paralel
  \item \textbf{Hardware Acceleration}: Gunakan AES-NI instruction
  \item \textbf{Constant-time}: Hindari branching berdasarkan data
\end{itemize}

\begin{lstlisting}[style=cstyle, caption={Optimasi dengan Tabel Lookup}]
// Pre-computed tables for rounds 1-9
static uint32_t te0[256], te1[256], te2[256], te3[256];
static uint32_t te4[256]; // For inverse

void init_tables() {
    for (int i = 0; i < 256; i++) {
        uint8_t s = sbox[i];
        te0[i] = (gmul(0x02, s) << 24) | (s << 16) | (s << 8) | gmul(0x03, s);
        te1[i] = (te0[i] << 8) | (te0[i] >> 24);
        te2[i] = (te1[i] << 8) | (te1[i] >> 24);
        te3[i] = (te2[i] << 8) | (te2[i] >> 24);
    }
}

// Optimized round (combines SubBytes, ShiftRows, MixColumns)
void aes_round_optimized(uint32_t *state, uint32_t *round_key) {
    state[0] ^= te0[(state[0] >> 24) & 0xff] ^
             te1[(state[1] >> 16) & 0xff] ^
             te2[(state[2] >> 8) & 0xff] ^
             te3[state[3] & 0xff] ^
             round_key[0];
    // ... similar for state[1], state[2], state[3]
}
\end{lstlisting}

\subsection{Keamanan Implementasi}

Implementasi AES yang aman harus memperhatikan:

\begin{itemize}
  \item \textbf{Constant-time operations}: Hindari timing attacks
  \item \textbf{Side-channel resistance}: Lindungi dari power analysis
  \item \textbf{Memory safety}: Hindari buffer overflow
  \item \textbf{Key protection}: Lindungi kunci dari memory dump
\end{itemize}

\begin{catatan}
Implementasi di atas adalah untuk tujuan edukasi. Untuk produksi, gunakan library yang telah diaudit keamanannya seperti OpenSSL, libsodium, atau implementasi hardware dengan AES-NI.
\end{catatan}
